\chapter{A Formal Model of Safety Boundaries}
\label{ch:formal}

\section{Motivation for a Formal Model}
\label{sec:formal-mot}

The phrase ``safety boundary'' has recurred throughout this monograph. To make
the architecture auditable rather than merely described, we now give a
precise, mathematical account of what a boundary is and what it means for one
to hold.

\section{Definition}
\label{sec:formal-def}

\begin{definition}
A \emph{safety boundary} $B$ is a quadruple
\[
B = (I_B,\, T_B,\, O_B,\, \phi_B)
\]
where:
\begin{itemize}
  \item $I_B$ is the inbound contract --- a predicate on artifacts
        $I_B : \mathcal{A} \to \{0,1\}$.
  \item $T_B : \mathcal{A} \to \mathcal{A}$ is the transformation.
  \item $O_B$ is the outbound contract --- a predicate
        $O_B : \mathcal{A} \to \{0,1\}$.
  \item $\phi_B$ is the \emph{consistency certificate}:
        $\phi_B(a) \in \{0,1\}$ records whether the transformation was
        verified.
\end{itemize}
\end{definition}

\begin{definition}
Boundary $B$ \emph{holds} on artifact $a$ iff
\[
I_B(a) = 1 \;\Longrightarrow\; O_B(T_B(a)) = 1 \;\wedge\; \phi_B(a) = 1.
\]
\end{definition}

In words: if the input satisfied the inbound contract, then the output must
satisfy the outbound contract and the transformation must be certified. If
$I_B(a) = 0$ the boundary \emph{fails loudly}: it must emit a sealed receipt
recording the failure and must not emit a (false) success artifact.

\section{Composing Boundaries}
\label{sec:formal-compose}

A pipeline is a sequence of boundaries $B_1,\dots,B_{10}$ with the
compatibility condition
\[
O_{B_i} \subseteq I_{B_{i+1}} \quad\text{for all } i,
\]
so that the output contract of one layer is accepted by the next. This is the
formal version of the ``artifact flow'' arrows in \cref{fig:vscp}.

\begin{theorem}
If every boundary $B_i$ holds on the artifact it receives, then the final
artifact $a_{10} = T_{B_{10}}(\dots T_{B_1}(a_0)\dots)$ satisfies
$O_{B_{10}}$ and carries a chain of certificates
$\phi_{B_1},\dots,\phi_{B_{10}}$.
\end{theorem}
\begin{proof}
By induction on $i$. Base: $B_1$ holds, so $O_{B_1}(a_1)=1$ and
$\phi_{B_1}=1$. Inductive step: assume $O_{B_i}(a_i)=1$; by compatibility
$I_{B_{i+1}}(a_i)=1$; since $B_{i+1}$ holds,
$O_{B_{i+1}}(a_{i+1})=1$ and $\phi_{B_{i+1}}=1$. The chain of certificates
is the concatenation of the individual $\phi$'s, each sealable into the WORM
ledger (Layer~9).
\end{proof}

\section{The Receipt as a Proof Object}
\label{sec:formal-receipt}

Each certificate $\phi_{B_i}$ is, concretely, a receipt
(\cref{ch:worm}) containing:
\begin{itemize}
  \item a SHA-256 of the input artifact $a_{i-1}$,
  \item a SHA-256 of the output artifact $a_i$,
  \item the boundary id $i$,
  \item an Ed25519 signature over the above (the Plasma Gate).
\end{itemize}
Because the receipt is itself an artifact, it becomes the input contract
evidence for the next boundary's verification, closing the loop between the
logical model and the cryptographic one.

\section{Failure Semantics}
\label{sec:formal-fail}

A key design choice is that a violated inbound contract is \emph{not} a
silent pass. Formally, the boundary function is total but its ``success''
projection is partial:
\[
\text{success}(B,a) =
\begin{cases}
1 & I_B(a)=1 \wedge O_B(T_B(a))=1 \wedge \phi_B(a)=1,\\
0 & I_B(a)=0 \text{ (recorded as a failed receipt)},\\
\bot & \text{the build gate halts}.
\end{cases}
\]
The build gate (\cref{ch:build}) realizes the $\bot$ case: any single
boundary returning $0$ aborts the build, so an unverified artifact can never
propagate.

\section{Relation to Established Theory}
\label{sec:formal-rel}

This model is a specialization of \emph{refinement} (Abrial's events /
Hoare's contracts): $I_B$ is the precondition, $O_B$ the postcondition, and
$\phi_B$ the verification witness. It is also an instance of a
\emph{layered trusted computing base}. The novelty is the explicit,
receipt-linked realization across a mixed-language symbolic-compute stack.
